\documentclass[11pt,letterpaper]{article}

% ------------------------------------------------------------
%  Packages
% ------------------------------------------------------------
\usepackage[utf8]{inputenc}
\usepackage[margin=1in]{geometry}
\usepackage{amsmath,amssymb,amsthm,mathtools}
\usepackage{booktabs}
\usepackage{enumitem}
\usepackage{titlesec}
\usepackage{xcolor}
\usepackage{etoolbox}
\usepackage{microtype}
\usepackage[skins,breakable]{tcolorbox}
\usepackage[parfill]{parskip}
\usepackage[hidelinks]{hyperref}

\hypersetup{
  pdftitle  = {MATH 225 -- Applied Linear Algebra II},
  pdfauthor = {Obrian Buchanan}
}

% ------------------------------------------------------------
%  Colours (clean, airy, modern)
% ------------------------------------------------------------
\definecolor{ink}{HTML}{1E293B}        % slate-900 (crisp modern charcoal)
\definecolor{textsub}{HTML}{64748B}    % slate-500 (subtle secondary text)

% Light, translucent-style background tints (96-98% lightness)
\definecolor{defshade}{HTML}{F8FAFC}   % slate-50 (ultra-light cool frost)
\definecolor{defborder}{HTML}{E2E8F0}  % slate-200 (hairline border)
\definecolor{thmshade}{HTML}{FBFBF9}   % warm alabaster tint
\definecolor{thmborder}{HTML}{E7E5E4}  % stone-200 (hairline border)

% ------------------------------------------------------------
%  Spacing
% ------------------------------------------------------------
\linespread{1.12}
\setlength{\parskip}{0.55em}
\setlength{\jot}{7pt}

\AtBeginDocument{%
  \setlength{\abovedisplayskip}{9pt plus 2pt minus 2pt}%
  \setlength{\belowdisplayskip}{9pt plus 2pt minus 2pt}%
  \setlength{\abovedisplayshortskip}{5pt plus 2pt}%
  \setlength{\belowdisplayshortskip}{7pt plus 2pt}%
}

\setlist{leftmargin=1.8em, topsep=0.5em, itemsep=0.5em, parsep=0.25em}

% Matrices: a little more room between rows and columns
\renewcommand{\arraystretch}{1.15}
\setlength{\arraycolsep}{4pt}

% Grey justification next to a derivation step, e.g. \by{V4}
\newcommand{\by}[1]{\text{\small\color{gray!80!black}#1}}

% ------------------------------------------------------------
%  Headings
% ------------------------------------------------------------
% Lecture title (section): number hidden in the body, shown in the TOC
\titleformat{\section}
  {\LARGE\bfseries\color{ink}}{}{0em}{}
\titlespacing*{\section}{0pt}{0pt}{2.2em}

% Subsection: numbered so each part is easy to spot
\titleformat{\subsection}
  {\Large\bfseries\color{ink}}{\thesubsection}{0.8em}{}
\titlespacing*{\subsection}{0pt}{2.6em}{1.0em}

% Subsubsection: smaller, unnumbered
\titleformat{\subsubsection}
  {\large\bfseries\color{ink}}{}{0em}{}
\titlespacing*{\subsubsection}{0pt}{1.8em}{0.6em}

\setcounter{tocdepth}{2}

% ------------------------------------------------------------
%  Numbering: one shared counter per lecture (Def 1.1, Thm 1.2, Ex 1.3, ...)
% ------------------------------------------------------------
\newcounter{note}[section]
\renewcommand{\thenote}{\thesection.\arabic{note}}

% ------------------------------------------------------------
%  Box styles (Modern frosted card aesthetic)
% ------------------------------------------------------------
\tcbset{
  notebox/.style={
    enhanced,
    boxrule=0.6pt,
    arc=4pt,
    left=14pt, right=14pt, top=10pt, bottom=10pt,
    before skip=1.5em, after skip=1.5em,
    detach title,
    fonttitle=\bfseries\small, coltitle=ink,
    before upper={\tcbtitle\par\vspace{0.4em}},
    opacityback=0.9
  },
  defbox/.style={
    notebox,
    colback=defshade,
    colframe=defborder
  },
  thmbox/.style={
    notebox,
    colback=thmshade,
    colframe=thmborder
  }
}

% Builds the box title, e.g. "Definition 1.2 (Subspace)"
\newcommand{\notetitle}[2]{#1~\thenote\ifstrempty{#2}{}{\ (#2)}}

% Definition: frosted slate card
\NewDocumentEnvironment{definition}{O{}}
  {\refstepcounter{note}\begin{tcolorbox}[defbox, title={\notetitle{Definition}{#1}}]}
  {\end{tcolorbox}}

% Theorem-type results: frosted stone card
\NewDocumentEnvironment{theorem}{O{}}
  {\refstepcounter{note}\begin{tcolorbox}[thmbox, title={\notetitle{Theorem}{#1}}]}
  {\end{tcolorbox}}

\NewDocumentEnvironment{proposition}{O{}}
  {\refstepcounter{note}\begin{tcolorbox}[thmbox, title={\notetitle{Proposition}{#1}}]}
  {\end{tcolorbox}}

\NewDocumentEnvironment{lemma}{O{}}
  {\refstepcounter{note}\begin{tcolorbox}[thmbox, title={\notetitle{Lemma}{#1}}]}
  {\end{tcolorbox}}

\NewDocumentEnvironment{corollary}{O{}}
  {\refstepcounter{note}\begin{tcolorbox}[thmbox, title={\notetitle{Corollary}{#1}}]}
  {\end{tcolorbox}}

% Key formula / checklist / method summary
\NewDocumentEnvironment{formula}{O{Key Formula / Checklist}}
  {\begin{tcolorbox}[thmbox, title={#1}]}
  {\end{tcolorbox}}

% Example: no box, clear heading with space around it
\NewDocumentEnvironment{example}{O{}}
  {\par\addvspace{1.6em}\refstepcounter{note}%
   \noindent{\bfseries\color{ink}\notetitle{Example}{#1}}\par\nopagebreak}
  {\par\addvspace{1.6em}}

% Solution: bold label, no end box
\NewDocumentEnvironment{solution}{}
  {\par\addvspace{0.4em}\noindent{\bfseries\color{ink}Solution.}\ }
  {\par}

% Intuition: plain paragraph with a coloured label
\NewDocumentEnvironment{intuition}{}
  {\par\addvspace{0.8em}\noindent{\bfseries\color{ink}Intuition.}\ }
  {\par\addvspace{0.8em}}

% Remark / note: plain paragraph with a label
\NewDocumentEnvironment{remark}{}
  {\par\addvspace{0.6em}\noindent{\bfseries\color{ink}Remark.}\ }
  {\par\addvspace{0.6em}}

% Proofs: space before and after
\AtBeginEnvironment{proof}{\vspace{0.3em}}
\AtEndEnvironment{proof}{\vspace{0.3em}}

% ============================================================
%  Document Body
% ============================================================
\begin{document}

% ------------------------------------------------------------
%  Title page
% ------------------------------------------------------------
\hypersetup{pageanchor=false}
\begin{titlepage}
    \vspace*{\fill}
    \begin{center}
        {\Huge\bfseries\color{ink} MATH 225\par}
        \vspace{0.4cm}
        {\LARGE\bfseries\color{ink} Applied Linear Algebra II\par}
        \vspace{1.2cm}
        {\Large Course Notes\par}
        \vspace{0.6cm}
        {\large Obrian Buchanan\par}
        \vspace{0.3cm}
        {\large University of Waterloo\par}
        \vspace{0.3cm}
        {\large Fall 2026\par}
    \end{center}
    \vspace*{\fill}
\end{titlepage}
\hypersetup{pageanchor=true}

% ------------------------------------------------------------
%  Table of contents
% ------------------------------------------------------------
{\setlength{\parskip}{0.3em}\tableofcontents}
\newpage

% ------------------------------------------------------------
%  Lectures
% ------------------------------------------------------------
\section{Lecture 1: General Vector Spaces and Subspaces}

\begin{intuition}
In previous courses, a ``vector'' was strictly an arrow in space (with magnitude and direction) or a column of numbers. In Applied Linear Algebra II, we generalize this concept entirely. A vector space is simply a mathematical sandbox. It is any collection of objects (whether they are matrices, polynomials, functions, or coordinates) that strictly obey a specific set of 10 rules. If they obey these rules, they are officially ``vectors,'' and all the powerful theorems of linear algebra apply to them.
\end{intuition}

The first five axioms govern how addition works internally, and the next five govern how scaling interacts with the vectors.

% =========================================================================
\subsection{Vector Spaces and the Ten Axioms}

\begin{definition}[Vector Space]
A \textbf{vector space over $\mathbb{R}$} is a set $V$ together with two operations, \textbf{addition} $\vec{x} + \vec{y}$ and \textbf{scalar multiplication} $c\vec{x}$, such that for all $\vec{x}, \vec{y}, \vec{z} \in V$ and all $c, d \in \mathbb{R}$ the following ten axioms hold.

\medskip
\textbf{Vector Addition Axioms}
\begin{enumerate}[label=\textbf{V\arabic*}, leftmargin=2.6em]
    \item \textbf{Closure under Addition:} $\vec{x} + \vec{y} \in V$
    \item \textbf{Commutativity:} $\vec{x} + \vec{y} = \vec{y} + \vec{x}$
    \item \textbf{Associativity:} $(\vec{x} + \vec{y}) + \vec{z} = \vec{x} + (\vec{y} + \vec{z})$
    \item \textbf{Additive Identity:} there is a $\vec{0} \in V$ with $\vec{x} + \vec{0} = \vec{x}$ for all $\vec{x} \in V$
    \item \textbf{Additive Inverse:} for each $\vec{x} \in V$ there is a $-\vec{x} \in V$ with $\vec{x} + (-\vec{x}) = \vec{0}$
\end{enumerate}

\medskip
\textbf{Scalar Multiplication Axioms}
\begin{enumerate}[label=\textbf{V\arabic*}, leftmargin=2.6em]
    \setcounter{enumi}{5}
    \item \textbf{Closure under Scalar Multiplication:} $c\vec{x} \in V$
    \item \textbf{Vector Distributivity:} $c(\vec{x} + \vec{y}) = c\vec{x} + c\vec{y}$
    \item \textbf{Scalar Distributivity:} $(c + d)\vec{x} = c\vec{x} + d\vec{x}$
    \item \textbf{Associativity:} $(cd)\vec{x} = c(d\vec{x})$
    \item \textbf{Multiplicative Identity:} $1\vec{x} = \vec{x}$
\end{enumerate}
\end{definition}

\begin{remark}
A set on its own is \emph{not} a vector space. We first have to say what ``addition'' and ``scalar multiplication'' mean on $V$, and only then check the ten axioms. The same set can be a vector space under one pair of operations and fail under another.
\end{remark}

\begin{example}[Checking All Ten Axioms]
Let $V = \mathbb{R}_{>0}$ be the set of positive real numbers, with ``addition'' and ``scalar multiplication'' defined by
\[
x \oplus y = xy, \qquad c \odot x = x^{c}.
\]
Show that $V$ is a vector space.

\begin{solution}
We go through the axioms one at a time. Throughout, $x, y, z > 0$ and $c, d \in \mathbb{R}$.

\medskip
\begin{center}
\renewcommand{\arraystretch}{1.3}
\begin{tabular}{@{}l l l@{}}
\textbf{V1} & $x \oplus y = xy > 0$ & so $x \oplus y \in V$ \\
\textbf{V2} & $x \oplus y = xy = yx = y \oplus x$ & \\
\textbf{V3} & $(x \oplus y) \oplus z = (xy)z = x(yz) = x \oplus (y \oplus z)$ & \\
\textbf{V4} & $x \oplus 1 = x \cdot 1 = x$ & so the zero vector is $1$ \\
\textbf{V5} & $x \oplus \tfrac{1}{x} = 1$ & so $-x = \tfrac{1}{x} \in V$ \\
\textbf{V6} & $c \odot x = x^{c} > 0$ & so $c \odot x \in V$ \\
\textbf{V7} & $c \odot (x \oplus y) = (xy)^{c} = x^{c}y^{c} = (c \odot x) \oplus (c \odot y)$ & \\
\textbf{V8} & $(c + d) \odot x = x^{c+d} = x^{c}x^{d} = (c \odot x) \oplus (d \odot x)$ & \\
\textbf{V9} & $(cd) \odot x = x^{cd} = (x^{d})^{c} = c \odot (d \odot x)$ & \\
\textbf{V10} & $1 \odot x = x^{1} = x$ & 
\end{tabular}
\end{center}
\medskip

All ten axioms hold, so $V$ is a vector space. Notice that the ``zero vector'' here is the number $1$, and the ``negative'' of $x$ is $\tfrac{1}{x}$. The zero vector is just whatever element acts as the additive identity, it does not have to be the number $0$.
\end{solution}
\end{example}

\subsubsection{Showing a Set Is Not a Vector Space}

To show something \emph{is} a vector space, all ten axioms have to hold. To show something is \emph{not} a vector space, we only need \textbf{one} axiom to fail, and we prove it fails with a specific counterexample (actual numbers).

\begin{example}[The Integers]
Let $V = \mathbb{Z}$ with the usual addition and multiplication. Is $V$ a vector space over $\mathbb{R}$?

\begin{solution}
No. Closure under scalar multiplication (V6) fails. Take $c = \tfrac{1}{2}$ and $x = 1 \in \mathbb{Z}$:
\[
cx = \tfrac{1}{2} \notin \mathbb{Z}.
\]
Scaling a vector in $V$ produced something outside $V$, so $\mathbb{Z}$ is not a vector space over $\mathbb{R}$.
\end{solution}
\end{example}

\begin{example}[A Strange Scalar Multiplication]
Let $V = \mathbb{R}^2$ with the usual addition, but with scalar multiplication defined by
\[
c \odot \begin{bmatrix} x \\ y \end{bmatrix} = \begin{bmatrix} cx \\ 0 \end{bmatrix}.
\]
Is $V$ a vector space?

\begin{solution}
No. The multiplicative identity axiom (V10) fails:
\[
1 \odot \begin{bmatrix} 1 \\ 1 \end{bmatrix} = \begin{bmatrix} 1 \\ 0 \end{bmatrix} \neq \begin{bmatrix} 1 \\ 1 \end{bmatrix}.
\]
The set $\mathbb{R}^2$ is fine, but with this operation it is not a vector space.
\end{solution}
\end{example}

% =========================================================================
\subsection{The Standard Vector Spaces}

The elements of a vector space $V$ are called \textbf{vectors}. The elements of the field (which for this course is always $\mathbb{R}$) are called \textbf{scalars}. These are the four spaces we will use over and over.

\begin{definition}[Euclidean Space $\mathbb{R}^n$]
\[
\mathbb{R}^n = \left\{ \begin{bmatrix} x_1 \\ \vdots \\ x_n \end{bmatrix} \;\middle|\; x_1, \dots, x_n \in \mathbb{R} \right\}
\]
with addition and scalar multiplication done component-wise:
\[
\begin{bmatrix} x_1 \\ \vdots \\ x_n \end{bmatrix} + \begin{bmatrix} y_1 \\ \vdots \\ y_n \end{bmatrix} = \begin{bmatrix} x_1 + y_1 \\ \vdots \\ x_n + y_n \end{bmatrix},
\qquad
c\begin{bmatrix} x_1 \\ \vdots \\ x_n \end{bmatrix} = \begin{bmatrix} cx_1 \\ \vdots \\ cx_n \end{bmatrix}.
\]
The zero vector is the column of $n$ zeros.
\end{definition}

\begin{example}[Computing in $\mathbb{R}^3$]
Let
\[
\vec{x} = \begin{bmatrix} 1 \\ -2 \\ 3 \end{bmatrix}, \qquad \vec{y} = \begin{bmatrix} 4 \\ 0 \\ -1 \end{bmatrix}.
\]
Compute $\vec{x} + \vec{y}$ and $2\vec{x} - \vec{y}$.

\begin{solution}
Working one entry at a time:
\[
\vec{x} + \vec{y} = \begin{bmatrix} 1 + 4 \\ -2 + 0 \\ 3 - 1 \end{bmatrix} = \begin{bmatrix} 5 \\ -2 \\ 2 \end{bmatrix},
\qquad
2\vec{x} - \vec{y} = \begin{bmatrix} 2 - 4 \\ -4 - 0 \\ 6 + 1 \end{bmatrix} = \begin{bmatrix} -2 \\ -4 \\ 7 \end{bmatrix}.
\]
\end{solution}
\end{example}

\begin{definition}[Matrix Space $M_{m \times n}(\mathbb{R})$]
$M_{m \times n}(\mathbb{R})$ is the set of all $m \times n$ matrices with real entries, with addition and scalar multiplication done entry by entry:
\[
(A + B)_{ij} = (A)_{ij} + (B)_{ij}, \qquad (cA)_{ij} = c\,(A)_{ij}.
\]
The zero vector is the zero matrix $O_{m \times n}$, consisting entirely of zeros.
\end{definition}

\begin{example}[Computing in $M_{2 \times 2}(\mathbb{R})$]
Let
\[
A = \begin{bmatrix} 1 & 2 \\ 0 & -1 \end{bmatrix}, \qquad B = \begin{bmatrix} 3 & -1 \\ 4 & 2 \end{bmatrix}.
\]
Compute $A + B$ and $A - 2B$, and write down the additive inverse of $A$.

\begin{solution}
Entry by entry:
\[
A + B = \begin{bmatrix} 4 & 1 \\ 4 & 1 \end{bmatrix},
\qquad
A - 2B = \begin{bmatrix} 1 - 6 & 2 + 2 \\ 0 - 8 & -1 - 4 \end{bmatrix} = \begin{bmatrix} -5 & 4 \\ -8 & -5 \end{bmatrix}.
\]
The additive inverse flips every sign, since then $A + (-A) = O_{2 \times 2}$:
\[
-A = \begin{bmatrix} -1 & -2 \\ 0 & 1 \end{bmatrix}.
\]
\end{solution}
\end{example}

\begin{definition}[Polynomial Space $P_n(\mathbb{R})$]
$P_n(\mathbb{R})$ is the set of all polynomials of degree \emph{at most} $n$ with real coefficients:
\[
P_n(\mathbb{R}) = \{ a_0 + a_1x + \dots + a_nx^n \mid a_0, a_1, \dots, a_n \in \mathbb{R} \},
\]
with the operations
\[
(p + q)(x) = p(x) + q(x), \qquad (cp)(x) = c\,p(x).
\]
In terms of coefficients, we add and scale matching coefficients. The zero vector is the zero polynomial $z(x) = 0 + 0x + \dots + 0x^n$.
\end{definition}

\begin{example}[Computing in $P_2(\mathbb{R})$]
Let $p(x) = 1 + 2x - x^2$ and $q(x) = 3 - x + 4x^2$. Compute $p + q$ and $2p - q$. Then explain why $P_2(\mathbb{R})$ is defined as degree \emph{at most} 2, not \emph{exactly} 2.

\begin{solution}
Combining matching coefficients:
\begin{align*}
p + q  &= (1 + 3) + (2 - 1)x + (-1 + 4)x^2 = 4 + x + 3x^2, \\
2p - q &= (2 + 4x - 2x^2) - (3 - x + 4x^2) = -1 + 5x - 6x^2.
\end{align*}
Now take $r(x) = 1 + x^2$ and $s(x) = 2 - x^2$. Both have degree exactly 2, but
\[
r + s = 3,
\]
which has degree 0. If we only allowed degree exactly 2, the set would not be closed under addition (V1), and it would not contain the zero polynomial at all. That is why we use ``degree at most $n$.''
\end{solution}
\end{example}

\begin{definition}[Function Space $\mathcal{F}(\mathbb{R})$]
$\mathcal{F}(\mathbb{R})$ is the set of all functions $f : \mathbb{R} \to \mathbb{R}$, with
\[
(f + g)(x) = f(x) + g(x), \qquad (cf)(x) = c\,f(x).
\]
The zero vector is the function that outputs $0$ for every $x$, and $-f$ is the function $(-f)(x) = -f(x)$.
\end{definition}

\begin{example}[Computing in $\mathcal{F}(\mathbb{R})$]
Let $f(x) = \sin x$ and $g(x) = x^2$. Describe $f + g$ and $3f$, and evaluate both at $x = \pi$.

\begin{solution}
By the definitions of the operations,
\[
(f + g)(x) = \sin x + x^2, \qquad (3f)(x) = 3\sin x.
\]
At $x = \pi$:
\[
(f + g)(\pi) = 0 + \pi^2 = \pi^2, \qquad (3f)(\pi) = 3 \cdot 0 = 0.
\]
The key idea is that each ``vector'' here is an entire function, and adding two vectors means building a new function point by point.
\end{solution}
\end{example}

% =========================================================================
\subsection{Fundamental Properties Derived from the Axioms}\label{sec:L1-props}

\begin{intuition}
Because we have redefined ``vectors'' to be anything from matrices to functions, we cannot blindly assume that basic arithmetic works the same way it does with regular numbers. We must formally prove that things like $0 \cdot \vec{x} = \vec{0}$ are guaranteed strictly by the 10 axioms, without relying on our intuition.
\end{intuition}

\begin{theorem}[Uniqueness of the Zero Vector and Inverses]\label{thm:L1-unique}
Let $V$ be a vector space over $\mathbb{R}$. Then:
\begin{enumerate}
    \item the zero vector $\vec{0}$ in $V$ is unique, and
    \item for each $\vec{x} \in V$, the additive inverse $-\vec{x}$ is unique.
\end{enumerate}
\end{theorem}

\begin{proof}
\textbf{Proof of (1).} Suppose $\vec{0}$ and $\vec{0}'$ both satisfy V4. Then
\begin{alignat*}{2}
\vec{0} &= \vec{0} + \vec{0}'  &\qquad& \text{($\vec{0}'$ is an identity)} \\
        &= \vec{0}' + \vec{0}  &&       \text{(by V2)} \\
        &= \vec{0}'            &&       \text{($\vec{0}$ is an identity)}
\end{alignat*}
So any two zero vectors are equal.

\medskip
\textbf{Proof of (2).} Suppose $\vec{y}$ and $\vec{z}$ are both inverses of $\vec{x}$, so $\vec{x} + \vec{y} = \vec{0}$ and $\vec{x} + \vec{z} = \vec{0}$. Then
\begin{alignat*}{2}
\vec{y} &= \vec{y} + \vec{0}                &\qquad& \text{(by V4)} \\
        &= \vec{y} + (\vec{x} + \vec{z})    &&       \text{($\vec{z}$ is an inverse of $\vec{x}$)} \\
        &= (\vec{x} + \vec{y}) + \vec{z}    &&       \text{(by V3, V2)} \\
        &= \vec{0} + \vec{z} = \vec{z}      &&       \text{($\vec{y}$ is an inverse, then V2, V4)}
\end{alignat*}
So the inverse of $\vec{x}$ is unique.
\end{proof}

\begin{theorem}[Elementary Arithmetic Properties of Vector Spaces]\label{thm:L1-arith}
Let $V$ be a vector space over $\mathbb{R}$. For all $\vec{x} \in V$ and all $t \in \mathbb{R}$:
\begin{enumerate}
    \item $0\vec{x} = \vec{0}$
    \item $t\vec{0} = \vec{0}$
    \item $(-1)\vec{x} = -\vec{x}$
    \item if $t\vec{x} = \vec{0}$, then either $t = 0$ or $\vec{x} = \vec{0}$.
\end{enumerate}
\end{theorem}

\begin{proof}
\textbf{Proof of (1).} In the real numbers, $0 + 0 = 0$, so by scalar distributivity (V8)
\[
0\vec{x} = (0 + 0)\vec{x} = 0\vec{x} + 0\vec{x}.
\]
Since $0\vec{x} \in V$ by closure (V6), its additive inverse $-(0\vec{x})$ exists in $V$ (V5). Adding this inverse to both sides:
\begin{alignat*}{2}
0\vec{x} + \bigl(-(0\vec{x})\bigr) &= \bigl(0\vec{x} + 0\vec{x}\bigr) + \bigl(-(0\vec{x})\bigr) && \\
\vec{0} &= 0\vec{x} + \bigl(0\vec{x} + (-(0\vec{x}))\bigr) &\qquad& \text{(by V5 on left, V3 on right)} \\
\vec{0} &= 0\vec{x} + \vec{0}                              &&       \text{(by V5)} \\
\vec{0} &= 0\vec{x}                                        &&       \text{(by V4)}
\end{alignat*}

\textbf{Proof of (2).} By the additive identity axiom, $\vec{0} + \vec{0} = \vec{0}$. Multiplying by $t$ and using vector distributivity (V7):
\[
t\vec{0} = t(\vec{0} + \vec{0}) = t\vec{0} + t\vec{0}.
\]
This is the same situation as in (1). Adding $-(t\vec{0})$ to both sides and repeating the same three steps (V5, V3, V4) gives $\vec{0} = t\vec{0}$.

\medskip
\textbf{Proof of (3).} We show that $(-1)\vec{x}$ behaves exactly as the additive inverse of $\vec{x}$:
\begin{alignat*}{2}
\vec{x} + (-1)\vec{x} &= 1\vec{x} + (-1)\vec{x}        &\qquad& \text{(by V10)} \\
                      &= \bigl(1 + (-1)\bigr)\vec{x}   &&       \text{(by V8)} \\
                      &= 0\vec{x} = \vec{0}            &&       \text{(by part (1))}
\end{alignat*}
So $(-1)\vec{x}$ is an inverse of $\vec{x}$. By uniqueness of inverses (Theorem~\ref{thm:L1-unique}), $(-1)\vec{x} = -\vec{x}$.

\medskip
\textbf{Proof of (4).} Suppose $t\vec{x} = \vec{0}$. If $t = 0$, the claim holds. If $t \neq 0$, then $t^{-1} \in \mathbb{R}$ exists, and multiplying both sides by $t^{-1}$:
\begin{alignat*}{2}
t^{-1}(t\vec{x}) &= t^{-1}\vec{0} && \\
(t^{-1}t)\vec{x} &= \vec{0}       &\qquad& \text{(by V9 on left, part (2) on right)} \\
\vec{x}          &= \vec{0}       &&       \text{(by $t^{-1}t = 1$, then V10)}
\end{alignat*}
So whenever $t \neq 0$, we are forced to have $\vec{x} = \vec{0}$.
\end{proof}

\begin{example}[Using the Properties]
Let $V$ be a vector space, $\vec{x} \in V$ and $t \in \mathbb{R}$. Show that $(-t)\vec{x} = -(t\vec{x})$.

\begin{solution}
Using the theorem above instead of going back to the axioms:
\begin{alignat*}{2}
(-t)\vec{x} &= \bigl((-1)t\bigr)\vec{x} &\qquad& \text{(arithmetic in $\mathbb{R}$)} \\
            &= (-1)(t\vec{x})           &&       \text{(by V9)} \\
            &= -(t\vec{x})               &&       \text{(by part (3), applied to $t\vec{x}$)}
\end{alignat*}
\end{solution}
\end{example}

% =========================================================================
\subsection{Subspaces and the Subspace Test}

\begin{intuition}
A subspace is simply a smaller, self-contained vector space living inside a larger one (for example, a 2D plane living inside a 3D space). Instead of verifying all 10 axioms to prove a subset is a vector space, we can take a shortcut. Since the subset inherits properties like commutativity and distributivity from the parent space, we only need to check if the subset is a ``closed system.'' If it contains the origin, and doesn't spill out of its boundaries when we add or scale elements, it is a subspace.
\end{intuition}

\begin{definition}[Subspace]
Let $V$ be a vector space over $\mathbb{R}$. A subset $U \subseteq V$ is called a \textbf{subspace} of $V$ if $U$ is itself a vector space over $\mathbb{R}$ under the addition and scalar multiplication operations inherited from $V$.
\end{definition}

\begin{example}[Using the Definition Directly]
In $\mathbb{R}^2$, consider the $x$-axis and a horizontal line shifted up by $1$:
\[
U = \left\{ \begin{bmatrix} a \\ 0 \end{bmatrix} \;\middle|\; a \in \mathbb{R} \right\},
\qquad
L = \left\{ \begin{bmatrix} a \\ 1 \end{bmatrix} \;\middle|\; a \in \mathbb{R} \right\}.
\]
Which of these is a subspace of $\mathbb{R}^2$?

\begin{solution}
For $U$: adding and scaling keeps the second entry at $0$,
\[
\begin{bmatrix} a \\ 0 \end{bmatrix} + \begin{bmatrix} b \\ 0 \end{bmatrix} = \begin{bmatrix} a + b \\ 0 \end{bmatrix} \in U,
\qquad
c\begin{bmatrix} a \\ 0 \end{bmatrix} = \begin{bmatrix} ca \\ 0 \end{bmatrix} \in U.
\]
The zero vector is in $U$ (take $a = 0$), and the inverse of a vector in $U$ is also in $U$ (replace $a$ by $-a$). Every other axiom is an equation that already holds in $\mathbb{R}^2$. So $U$ is a vector space on its own, and hence a subspace.

For $L$: adding two vectors of $L$ leaves $L$,
\[
\begin{bmatrix} 0 \\ 1 \end{bmatrix} + \begin{bmatrix} 0 \\ 1 \end{bmatrix} = \begin{bmatrix} 0 \\ 2 \end{bmatrix} \notin L,
\]
so V1 fails and $L$ is not a subspace.

Notice what we actually had to check for $U$: the zero vector and the two closure properties. Everything else came for free. That is exactly the shortcut below.
\end{solution}
\end{example}

\begin{theorem}[Subspace Test]
Let $V$ be a vector space over $\mathbb{R}$. A subset $U \subseteq V$ is a subspace of $V$ if and only if:
\begin{enumerate}
    \item $\vec{0} \in U$, where $\vec{0}$ is the zero vector of $V$ (in particular, $U \neq \emptyset$),
    \item for all $\vec{x}, \vec{y} \in U$, $\vec{x} + \vec{y} \in U$ \quad (\textbf{closure under addition}),
    \item for all $\alpha \in \mathbb{R}$ and $\vec{x} \in U$, $\alpha \vec{x} \in U$ \quad (\textbf{closure under scalar multiplication}).
\end{enumerate}
\end{theorem}

\begin{proof}
\textbf{($\Longrightarrow$)} If $U$ is a subspace, it satisfies all the axioms on its own. In particular it is closed under addition (V1) and scalar multiplication (V6), which gives (2) and (3). For (1), $U$ is non-empty, so pick any $\vec{x} \in U$. By closure $0\vec{x} \in U$, and $0\vec{x} = \vec{0}$ by Theorem~\ref{thm:L1-arith}(1). So $\vec{0} \in U$.

\medskip
\textbf{($\Longleftarrow$)} Assume (1), (2) and (3). We check the ten axioms for $U$:
\begin{itemize}
    \item \textbf{V1, V6:} these are exactly conditions (2) and (3).
    \item \textbf{V4:} $\vec{0} \in U$ by (1), and $\vec{x} + \vec{0} = \vec{x}$ already holds in $V$.
    \item \textbf{V5:} for $\vec{x} \in U$, condition (3) with $\alpha = -1$ gives $(-1)\vec{x} \in U$, and $(-1)\vec{x} = -\vec{x}$ by Theorem~\ref{thm:L1-arith}(3).
    \item \textbf{V2, V3, V7 to V10:} these are equations that hold for \emph{all} vectors of $V$, so they hold for the vectors in $U$.
\end{itemize}
Thus $U$ is a vector space, and so a subspace of $V$.
\end{proof}

\begin{remark}
Condition (1) is the quickest check, so always do it first. If $\vec{0} \notin U$, you are done: $U$ is not a subspace. But $\vec{0} \in U$ on its own proves nothing, you still need both closure conditions.
\end{remark}

\begin{remark}
Every vector space $V$ has two \textbf{trivial subspaces}: $\{\vec{0}\}$ and $V$ itself. Also, a subspace has to be a \emph{subset} of $V$, so $\mathbb{R}^2$ is not a subspace of $\mathbb{R}^3$ (its vectors have two entries, not three). More generally, $\mathbb{R}^k$ is not a subspace of $\mathbb{R}^n$ whenever $k \neq n$.
\end{remark}

\subsubsection{Geometric Picture}

In $\mathbb{R}^2$ and $\mathbb{R}^3$, subspaces are exactly the ``flat'' sets passing through the origin:
\begin{itemize}
    \item subspaces of $\mathbb{R}^2$: $\{\vec{0}\}$, lines through the origin, and $\mathbb{R}^2$;
    \item subspaces of $\mathbb{R}^3$: $\{\vec{0}\}$, lines through the origin, planes through the origin, and $\mathbb{R}^3$.
\end{itemize}
Anything curved (like a parabola), anything shifted off the origin (like the line $y = x + 1$), or anything with an edge (like a quadrant) fails.

% =========================================================================
\subsection{Worked Examples: Subspace Analysis}

\begin{formula}[How to Run the Subspace Test]
\textbf{To prove $U$ is a subspace:}
\begin{enumerate}
    \item Show the zero vector of $V$ satisfies the condition that defines $U$.
    \item Take two \emph{general} elements of $U$, write down what membership means, and show their sum satisfies the same condition.
    \item Take a \emph{general} element of $U$ and a general scalar $\alpha$, and show the scaled element satisfies the condition.
\end{enumerate}
\textbf{To prove $U$ is not a subspace:} find \emph{one} condition that fails, using specific numbers as a counterexample.
\end{formula}

\begin{example}[Polynomial with Root Constraint]
Let $V = P_2(\mathbb{R})$ and
\[
U = \{p(x) \in P_2(\mathbb{R}) \mid p(3) = 0\}.
\]
Prove that $U$ is a subspace of $P_2(\mathbb{R})$.

\begin{solution}
We test the three conditions of the Subspace Test.

\medskip
\textbf{Zero vector.} The zero polynomial $z(x) = 0 + 0x + 0x^2$ gives $z(3) = 0$, so $z \in U$.

\medskip
\textbf{Closure under addition.} Let $p, q \in U$, so $p(3) = 0$ and $q(3) = 0$. Then
\[
(p + q)(3) = p(3) + q(3) = 0 + 0 = 0,
\]
so $p + q \in U$.

\medskip
\textbf{Closure under scalar multiplication.} Let $p \in U$ and $\alpha \in \mathbb{R}$. Then
\[
(\alpha p)(3) = \alpha\, p(3) = \alpha \cdot 0 = 0,
\]
so $\alpha p \in U$.

All three conditions hold, so $U$ is a subspace of $P_2(\mathbb{R})$.
\end{solution}
\end{example}

\begin{example}[Traceless Matrices]
Let $V = M_{2\times 2}(\mathbb{R})$ and
\[
W = \{ A \in M_{2\times 2}(\mathbb{R}) \mid \operatorname{tr}(A) = 0 \},
\]
where the trace is the sum of the main diagonal:
\[
\operatorname{tr}\begin{bmatrix} a & b \\ c & d \end{bmatrix} = a + d.
\]
Show that $W$ is a subspace of $M_{2\times 2}(\mathbb{R})$.

\begin{solution}
\textbf{Zero vector.} $\operatorname{tr}(O_{2\times 2}) = 0 + 0 = 0$, so $O_{2\times 2} \in W$.

\medskip
\textbf{Closure under addition.} Let $A, B \in W$, with
\[
A = \begin{bmatrix} a_1 & b_1 \\ c_1 & d_1 \end{bmatrix}, \qquad
B = \begin{bmatrix} a_2 & b_2 \\ c_2 & d_2 \end{bmatrix},
\]
so $a_1 + d_1 = 0$ and $a_2 + d_2 = 0$. Then
\[
\operatorname{tr}(A + B) = (a_1 + a_2) + (d_1 + d_2) = (a_1 + d_1) + (a_2 + d_2) = 0,
\]
so $A + B \in W$.

\medskip
\textbf{Closure under scalar multiplication.} Let $A \in W$ and $\alpha \in \mathbb{R}$. Then
\[
\operatorname{tr}(\alpha A) = \alpha a_1 + \alpha d_1 = \alpha(a_1 + d_1) = 0,
\]
so $\alpha A \in W$.

Thus $W$ is a subspace of $M_{2\times 2}(\mathbb{R})$.
\end{solution}
\end{example}

\begin{example}[Matrices with Equal Off-Diagonal Entries]
Determine whether
\[
S = \left\{ \begin{bmatrix} a_1 & a_2 \\ a_3 & a_4 \end{bmatrix} \in M_{2\times 2}(\mathbb{R}) \;\middle|\; a_2 = a_3 \right\}
\]
is a subspace of $M_{2\times 2}(\mathbb{R})$.

\begin{solution}
\textbf{Zero vector.} The zero matrix has $a_2 = 0 = a_3$, so $O_{2\times 2} \in S$.

\medskip
\textbf{Closure under addition.} Let $A, B \in S$, so $a_2 = a_3$ and $b_2 = b_3$. Then
\[
A + B = \begin{bmatrix} a_1 + b_1 & a_2 + b_2 \\ a_3 + b_3 & a_4 + b_4 \end{bmatrix},
\]
and the off-diagonal entries are equal since $a_2 + b_2 = a_3 + b_3$. So $A + B \in S$.

\medskip
\textbf{Closure under scalar multiplication.} Let $A \in S$ and $\alpha \in \mathbb{R}$. Then
\[
\alpha A = \begin{bmatrix} \alpha a_1 & \alpha a_2 \\ \alpha a_3 & \alpha a_4 \end{bmatrix},
\]
and $\alpha a_2 = \alpha a_3$. So $\alpha A \in S$.

Therefore $S$ is a subspace of $M_{2\times 2}(\mathbb{R})$.
\end{solution}
\end{example}

\begin{example}[A Plane Through the Origin in $\mathbb{R}^3$]
Let
\[
U = \left\{ \begin{bmatrix} x \\ y \\ z \end{bmatrix} \in \mathbb{R}^3 \;\middle|\; x - 2y + 3z = 0 \right\}.
\]
Prove that $U$ is a subspace of $\mathbb{R}^3$.

\begin{solution}
\textbf{Zero vector.} $0 - 2(0) + 3(0) = 0$, so $\vec{0} \in U$.

\medskip
\textbf{Closure under addition.} Let
\[
\vec{u} = \begin{bmatrix} x_1 \\ y_1 \\ z_1 \end{bmatrix}, \qquad
\vec{v} = \begin{bmatrix} x_2 \\ y_2 \\ z_2 \end{bmatrix}
\]
be in $U$, so $x_1 - 2y_1 + 3z_1 = 0$ and $x_2 - 2y_2 + 3z_2 = 0$. The entries of $\vec{u} + \vec{v}$ satisfy
\begin{align*}
(x_1 + x_2) - 2(y_1 + y_2) + 3(z_1 + z_2)
&= (x_1 - 2y_1 + 3z_1) + (x_2 - 2y_2 + 3z_2) \\
&= 0 + 0 = 0,
\end{align*}
so $\vec{u} + \vec{v} \in U$.

\medskip
\textbf{Closure under scalar multiplication.} Let $\vec{u} \in U$ and $\alpha \in \mathbb{R}$. The entries of $\alpha\vec{u}$ satisfy
\[
(\alpha x_1) - 2(\alpha y_1) + 3(\alpha z_1) = \alpha(x_1 - 2y_1 + 3z_1) = \alpha \cdot 0 = 0,
\]
so $\alpha\vec{u} \in U$.

Since it passes the Subspace Test, the plane $U$ is a subspace of $\mathbb{R}^3$.
\end{solution}
\end{example}

\begin{example}[Affine Coefficient Failure]
Let $V = P_2(\mathbb{R})$ and
\[
S = \{ ax^2 + bx + c \in P_2(\mathbb{R}) \mid a - 2b + c = 3 \}.
\]
Determine whether $S$ is a subspace of $V$.

\begin{solution}
We always check the zero vector first. The zero polynomial has $a = b = c = 0$, and
\[
a - 2b + c = 0 \neq 3.
\]
So the zero polynomial is not in $S$, and condition (1) of the Subspace Test fails immediately. A subspace \emph{must} always contain the zero vector.

Geometrically, in the ``space'' of coefficients $(a, b, c)$, this is a plane that has been shifted off the origin. Thus $S$ is \textbf{not} a subspace of $P_2(\mathbb{R})$.
\end{solution}
\end{example}

\begin{example}[Contains Zero, Still Fails]
Determine whether
\[S = \left\{ \begin{bmatrix} x_1 \\ x_2 \\ x_3 \end{bmatrix} \in \mathbb{R}^3 \;\middle\vert{}\; x_1 = x_2 x_3 \right\}\]
is a subspace of $\mathbb{R}^3$.

\begin{solution}
The zero vector passes, since $0 = 0 \cdot 0$, so we look at closure. The condition multiplies two entries together, which is a red flag (it is not linear), so we look for a counterexample. Take
\[\vec{x} = \begin{bmatrix} 1 \\ 1 \\ 1 \end{bmatrix}, \qquad \vec{y} = \begin{bmatrix} 2 \\ 1 \\ 2 \end{bmatrix}, \qquad \vec{x} + \vec{y} = \begin{bmatrix} 3 \\ 2 \\ 3 \end{bmatrix}.\]
Both $\vec{x}$ and $\vec{y}$ are in $S$ since $1 = 1 \cdot 1$ and $2 = 1 \cdot 2$. But for the sum, $3 \neq 2 \cdot 3$, so $\vec{x} + \vec{y} \notin S$. Since $S$ is not closed under addition, it is \textbf{not} a subspace of $\mathbb{R}^3$.
\end{solution}
\end{example}

\begin{example}[The First Quadrant]
Let
\[Q = \left\{ \begin{bmatrix} x \\ y \end{bmatrix} \in \mathbb{R}^2 \;\middle\vert{}\; x \ge 0, \ y \ge 0 \right\}.\]
Is $Q$ a subspace of $\mathbb{R}^2$?

\begin{solution}
The zero vector is in $Q$, and $Q$ is even closed under addition (adding two non-negative numbers gives a non-negative number). But scalar multiplication fails:
\[(-1)\begin{bmatrix} 1 \\ 1 \end{bmatrix} = \begin{bmatrix} -1 \\ -1 \end{bmatrix} \notin Q.\]
So $Q$ is \textbf{not} a subspace of $\mathbb{R}^2$. This is why we need all three conditions: passing two of them is not enough.
\end{solution}
\end{example}
\newpage
\section{Lecture 2: Span and Linear Independence}

\begin{intuition}
If we treat vectors like directional building blocks, the ``span'' is the set of absolutely everywhere you can possibly reach by combining them. ``Linear independence'' means there is no redundancy in your building blocks; no vector in the set can be created by mixing the others. If a set spans a space and is linearly independent, it forms a perfect, highly efficient coordinate system for that space.
\end{intuition}

% =========================================================================
\subsection{Linear Combinations and Span}

\begin{definition}[Linear Combination]
Let $V$ be a vector space over $\mathbb{R}$ and let $S = \{\vec{v}_1, \dots, \vec{v}_k\} \subseteq V$. An expression of the form
\[
c_1\vec{v}_1 + c_2\vec{v}_2 + \dots + c_k\vec{v}_k, \qquad c_1, \dots, c_k \in \mathbb{R},
\]
is called a \textbf{linear combination} of the vectors in $S$. The scalars $c_i$ are called its \textbf{coefficients}.
\end{definition}

\begin{example}[Writing a Vector as a Linear Combination]
Write
\[
\vec{w} = \begin{bmatrix} -6 \\ 7 \\ 1 \end{bmatrix}
\quad \text{as a linear combination of} \quad
\vec{v}_1 = \begin{bmatrix} -2 \\ 2 \\ 1 \end{bmatrix}, \quad
\vec{v}_2 = \begin{bmatrix} 8 \\ -11 \\ 2 \end{bmatrix}.
\]

\begin{solution}
We want $c_1, c_2$ with $c_1\vec{v}_1 + c_2\vec{v}_2 = \vec{w}$. Comparing entries gives the system
\begin{align*}
-2c_1 + 8c_2 &= -6, \\
2c_1 - 11c_2 &= 7, \\
c_1 + 2c_2 &= 1.
\end{align*}
From the third equation, $c_1 = 1 - 2c_2$. Substituting into the first:
\[
-2(1 - 2c_2) + 8c_2 = -6 \implies 12c_2 = -4 \implies c_2 = -\tfrac{1}{3}, \qquad c_1 = 1 + \tfrac{2}{3} = \tfrac{5}{3}.
\]
We must check these also satisfy the second equation: $2\left(\tfrac{5}{3}\right) - 11\left(-\tfrac{1}{3}\right) = \tfrac{10}{3} + \tfrac{11}{3} = 7$. They do, so
\[
\vec{w} = \tfrac{5}{3}\,\vec{v}_1 - \tfrac{1}{3}\,\vec{v}_2.
\]
\end{solution}
\end{example}

\begin{definition}[Span]
Let $S = \{\vec{v}_1, \dots, \vec{v}_k\} \subseteq V$. The \textbf{span} of $S$ is the set of all possible linear combinations of the vectors in $S$:
\[
\operatorname{Span}(S) = \{ c_1\vec{v}_1 + \dots + c_k\vec{v}_k \mid c_1, \dots, c_k \in \mathbb{R} \}.
\]
If $S = \emptyset$, we define $\operatorname{Span}(\emptyset) = \{\vec{0}\}$.

\medskip
If $\operatorname{Span}(S) = W$, we say that $S$ \textbf{spans} $W$, or that $S$ is a \textbf{spanning set} for $W$.
\end{definition}

\begin{example}[Elements In and Out of a Span]
Let $S = \{x, \ 1 + x^2\} \subseteq P_2(\mathbb{R})$. Give some elements of $\operatorname{Span}(S)$, and show that $p(x) = 1$ is \emph{not} in $\operatorname{Span}(S)$.

\begin{solution}
Any choice of coefficients gives an element of the span:
\begin{alignat*}{2}
0 &= 0 \cdot x + 0 \cdot (1 + x^2) &\qquad& \by{$c_1 = 0,\ c_2 = 0$} \\
x &= 1 \cdot x + 0 \cdot (1 + x^2) && \by{$c_1 = 1,\ c_2 = 0$} \\
-2 + 3x - 2x^2 &= 3 \cdot x - 2 \cdot (1 + x^2) && \by{$c_1 = 3,\ c_2 = -2$}
\end{alignat*}
Now suppose $1 \in \operatorname{Span}(S)$. Then for some $c_1, c_2$,
\[
1 = c_1 x + c_2(1 + x^2) = c_2 + c_1 x + c_2 x^2.
\]
Matching coefficients, the constant term needs $c_2 = 1$ but the $x^2$ term needs $c_2 = 0$. That is impossible, so $1 \notin \operatorname{Span}(S)$.
\end{solution}
\end{example}

\begin{example}[Span Membership in $P_2(\mathbb{R})$]
Let $S = \{1, \ 1 - x - x^2, \ 1 - 2x - 2x^2\}$. Is $p(x) = -3 - x + 2x^2$ in $\operatorname{Span}(S)$?

\begin{solution}
Expand a general linear combination and group by powers of $x$:
\[
c_1(1) + c_2(1 - x - x^2) + c_3(1 - 2x - 2x^2) = (c_1 + c_2 + c_3) + (-c_2 - 2c_3)x + (-c_2 - 2c_3)x^2.
\]
Matching this with $-3 - x + 2x^2$, the $x$ and $x^2$ coefficients give
\[
-c_2 - 2c_3 = -1 \qquad \text{and} \qquad -c_2 - 2c_3 = 2.
\]
The same expression cannot equal both $-1$ and $2$, so the system is inconsistent. Therefore $p \notin \operatorname{Span}(S)$.
\end{solution}
\end{example}

% =========================================================================
\subsection{Span Is a Subspace}

\begin{theorem}[Span Is a Subspace]\label{thm:L2-span-subspace}
Let $V$ be a vector space over $\mathbb{R}$ and let $S = \{\vec{v}_1, \dots, \vec{v}_k\} \subseteq V$. Then $\operatorname{Span}(S)$ is a subspace of $V$.
\end{theorem}

\begin{proof}
We verify the three conditions of the Subspace Test.

\textbf{Zero vector.} Choosing $c_1 = \dots = c_k = 0$ gives $\vec{0} = 0\vec{v}_1 + \dots + 0\vec{v}_k \in \operatorname{Span}(S)$.

\textbf{Closure under addition.} Let $\vec{x}, \vec{y} \in \operatorname{Span}(S)$. Then there are scalars $c_i, d_i \in \mathbb{R}$ with
\[
\vec{x} = \sum_{i=1}^k c_i\vec{v}_i, \qquad \vec{y} = \sum_{i=1}^k d_i\vec{v}_i.
\]
Adding and using distributivity,
\[
\vec{x} + \vec{y} = \sum_{i=1}^k c_i\vec{v}_i + \sum_{i=1}^k d_i\vec{v}_i = \sum_{i=1}^k (c_i + d_i)\vec{v}_i.
\]
Each $c_i + d_i \in \mathbb{R}$, so $\vec{x} + \vec{y} \in \operatorname{Span}(S)$.

\textbf{Closure under scalar multiplication.} For any $\alpha \in \mathbb{R}$,
\[
\alpha \vec{x} = \alpha \sum_{i=1}^k c_i\vec{v}_i = \sum_{i=1}^k (\alpha c_i)\vec{v}_i.
\]
Each $\alpha c_i \in \mathbb{R}$, so $\alpha \vec{x} \in \operatorname{Span}(S)$.
\end{proof}

\begin{theorem}[Span Is the Smallest Subspace Containing $S$]
Let $S = \{\vec{v}_1, \dots, \vec{v}_k\} \subseteq V$. If $W$ is any subspace of $V$ that contains $\vec{v}_1, \dots, \vec{v}_k$, then
\[
\operatorname{Span}(S) \subseteq W.
\]
\end{theorem}

\begin{proof}
Take any $c_1\vec{v}_1 + \dots + c_k\vec{v}_k \in \operatorname{Span}(S)$. Each $\vec{v}_i \in W$, so each $c_i\vec{v}_i \in W$ by closure under scalar multiplication. Adding them one at a time, the whole sum is in $W$ by closure under addition. So every element of $\operatorname{Span}(S)$ is in $W$.
\end{proof}

\begin{remark}
Together, these two theorems say $\operatorname{Span}(S)$ is the \emph{minimal} subspace containing $S$: it is a subspace, and every other subspace containing $S$ is at least as big. This also gives a fast way to prove a set is a subspace: if you can write it as a span, you are done.
\end{remark}

\begin{example}[Proving a Subspace by Writing It as a Span]
Let $U = \{ p(x) \in P_2(\mathbb{R}) \mid p(1) = 0 \}$. Find a spanning set for $U$, and conclude that $U$ is a subspace.

\begin{solution}
Write $p(x) = a + bx + cx^2$. The condition $p(1) = 0$ means $a + b + c = 0$, so $a = -b - c$. Substituting,
\begin{align*}
p(x) &= (-b - c) + bx + cx^2 \\
     &= b(x - 1) + c(x^2 - 1).
\end{align*}
Every element of $U$ is a linear combination of $x - 1$ and $x^2 - 1$, and both of these satisfy $p(1) = 0$. So
\[
U = \operatorname{Span}\{x - 1, \ x^2 - 1\},
\]
which is a subspace of $P_2(\mathbb{R})$ by Theorem~\ref{thm:L2-span-subspace}.
\end{solution}
\end{example}

% =========================================================================
\subsection{Worked Examples: Proving a Set Spans a Space}

\begin{formula}[How to Check if $S$ Spans $V$]
To prove that a set of vectors $S$ spans an entire vector space $V$, you must show that an \emph{arbitrary} (generic) vector in $V$ can always be built using a linear combination of the vectors in $S$.
\begin{enumerate}
    \item Write a generic vector of $V$ with letters (e.g.\ $x, y$) and set it equal to $c_1\vec{v}_1 + \dots + c_k\vec{v}_k$.
    \item Turn this into a system of linear equations in the $c_i$.
    \item If the system is consistent for \emph{every} choice of the letters, $S$ spans $V$. If some choice makes it inconsistent, $S$ does not span $V$.
\end{enumerate}
\end{formula}

\begin{example}[Spanning $\mathbb{R}^2$]
Determine whether
\[
S = \left\{ \begin{bmatrix} 1 \\ 1 \end{bmatrix}, \begin{bmatrix} 1 \\ -1 \end{bmatrix} \right\}
\]
spans $\mathbb{R}^2$.

\begin{solution}
Let $(x, y)$ stand for the entries of an arbitrary vector in $\mathbb{R}^2$. We must determine if there always exist $c_1, c_2 \in \mathbb{R}$ with
\[
c_1 \begin{bmatrix} 1 \\ 1 \end{bmatrix} + c_2 \begin{bmatrix} 1 \\ -1 \end{bmatrix} = \begin{bmatrix} x \\ y \end{bmatrix}.
\]
This vector equation produces the system
\begin{alignat*}{2}
c_1 + c_2 &= x &\qquad& \by{Equation 1} \\
c_1 - c_2 &= y && \by{Equation 2}
\end{alignat*}
Adding the two equations eliminates $c_2$, and subtracting Equation 2 from Equation 1 eliminates $c_1$:
\[
2c_1 = x + y \implies c_1 = \frac{x + y}{2},
\qquad
2c_2 = x - y \implies c_2 = \frac{x - y}{2}.
\]
Because we can find explicit, real values for $c_1$ and $c_2$ regardless of what $x$ and $y$ are, every vector in $\mathbb{R}^2$ can be constructed. Therefore $S$ spans $\mathbb{R}^2$.
\end{solution}
\end{example}

\begin{example}[A Set That Does Not Span $\mathbb{R}^2$]
Determine whether
\[
S = \left\{ \begin{bmatrix} 1 \\ 2 \end{bmatrix}, \begin{bmatrix} 2 \\ 4 \end{bmatrix} \right\}
\]
spans $\mathbb{R}^2$.

\begin{solution}
Set
\[
c_1 \begin{bmatrix} 1 \\ 2 \end{bmatrix} + c_2 \begin{bmatrix} 2 \\ 4 \end{bmatrix} = \begin{bmatrix} x \\ y \end{bmatrix}
\]
and row reduce the augmented matrix:
\[
\left[\begin{array}{cc|c} 1 & 2 & x \\ 2 & 4 & y \end{array}\right]
\xrightarrow{R_2 - 2R_1}
\left[\begin{array}{cc|c} 1 & 2 & x \\ 0 & 0 & y - 2x \end{array}\right].
\]
The system is consistent only when $y - 2x = 0$. So, for example, the vector with $x = 1$, $y = 0$ cannot be reached, since $0 - 2(1) \neq 0$. Therefore $S$ does \textbf{not} span $\mathbb{R}^2$.

Geometrically, the second vector is just $2$ times the first, so both point along the same line $y = 2x$, and that line is the whole span.
\end{solution}
\end{example}

% =========================================================================
\subsection{Linear Independence}

\begin{definition}[Linear Dependence and Independence]
Let $V$ be a vector space over $\mathbb{R}$ and let $S = \{\vec{v}_1, \dots, \vec{v}_k\} \subseteq V$.
\begin{itemize}
    \item $S$ is \textbf{linearly dependent} if there exist scalars $c_1, \dots, c_k \in \mathbb{R}$, \textbf{not all zero}, such that
    \[
    c_1\vec{v}_1 + c_2\vec{v}_2 + \dots + c_k\vec{v}_k = \vec{0}.
    \]
    \item $S$ is \textbf{linearly independent} if the only solution to
    \[
    c_1\vec{v}_1 + c_2\vec{v}_2 + \dots + c_k\vec{v}_k = \vec{0}
    \]
    is the \textbf{trivial solution} $c_1 = c_2 = \dots = c_k = 0$.
\end{itemize}
\end{definition}

\begin{formula}[How to Test for Linear Independence]
To test for linear independence, we always set a generic linear combination of the set equal to the zero vector. We then solve for the coefficients $c_i$.
\begin{itemize}
    \item If we get a unique solution ($c_1 = c_2 = \dots = 0$), the set is \textbf{independent}.
    \item If we get free variables (meaning infinitely many non-zero solutions exist), the set is \textbf{dependent}, and one vector is just a combination of the others.
\end{itemize}
\end{formula}

\begin{example}[Independence in $\mathbb{R}^2$]
Determine whether
\[
S = \left\{ \begin{bmatrix} 2 \\ 3 \end{bmatrix}, \begin{bmatrix} -3 \\ 5 \end{bmatrix} \right\}
\]
is linearly independent.

\begin{solution}
Setting
\[
c_1 \begin{bmatrix} 2 \\ 3 \end{bmatrix} + c_2 \begin{bmatrix} -3 \\ 5 \end{bmatrix} = \begin{bmatrix} 0 \\ 0 \end{bmatrix}
\]
gives the system
\[
2c_1 - 3c_2 = 0, \qquad 3c_1 + 5c_2 = 0.
\]
From the first equation, $c_1 = \tfrac{3}{2}c_2$. Substituting into the second,
\[
\tfrac{9}{2}c_2 + 5c_2 = \tfrac{19}{2}c_2 = 0 \implies c_2 = 0 \implies c_1 = 0.
\]
Only the trivial solution exists, so $S$ is \textbf{linearly independent}.
\end{solution}
\end{example}

\begin{example}[Linear Dependence in $\mathbb{R}^3$]
Determine whether
\[
S = \left\{ \begin{bmatrix} 1 \\ 2 \\ 1 \end{bmatrix}, \begin{bmatrix} 2 \\ 5 \\ 0 \end{bmatrix}, \begin{bmatrix} 3 \\ 7 \\ 1 \end{bmatrix} \right\}
\]
is linearly independent.

\begin{solution}
We set up the dependence equation
\[
c_1 \begin{bmatrix} 1 \\ 2 \\ 1 \end{bmatrix} + c_2 \begin{bmatrix} 2 \\ 5 \\ 0 \end{bmatrix} + c_3 \begin{bmatrix} 3 \\ 7 \\ 1 \end{bmatrix} = \begin{bmatrix} 0 \\ 0 \\ 0 \end{bmatrix}
\]
and row reduce its augmented matrix:
\begin{align*}
\left[\begin{array}{ccc|c}
1 & 2 & 3 & 0 \\
2 & 5 & 7 & 0 \\
1 & 0 & 1 & 0
\end{array}\right]
&\xrightarrow[R_3 - R_1]{R_2 - 2R_1}
\left[\begin{array}{ccc|c}
1 & 2 & 3 & 0 \\
0 & 1 & 1 & 0 \\
0 & -2 & -2 & 0
\end{array}\right] \\[0.6em]
&\xrightarrow{R_3 + 2R_2}
\left[\begin{array}{ccc|c}
1 & 2 & 3 & 0 \\
0 & 1 & 1 & 0 \\
0 & 0 & 0 & 0
\end{array}\right].
\end{align*}
Because the third column has no pivot, $c_3$ is a free variable, so there are infinitely many non-zero solutions. For instance, let $c_3 = -1$. Then:
\begin{alignat*}{2}
c_2 + c_3 = 0 &\implies c_2 = 1 &\qquad& \by{row 2} \\
c_1 + 2c_2 + 3c_3 = 0 &\implies c_1 + 2 - 3 = 0 \implies c_1 = 1 && \by{row 1}
\end{alignat*}
So
\[
1\vec{v}_1 + 1\vec{v}_2 - 1\vec{v}_3 = \vec{0},
\]
a non-trivial combination equal to zero. The set is \textbf{linearly dependent}. (Notice $\vec{v}_3 = \vec{v}_1 + \vec{v}_2$.)
\end{solution}
\end{example}

\begin{example}[Linear Independence in $P_2(\mathbb{R})$]
Determine whether $S = \{1 + x, \ x + x^2, \ 1 + x^2\}$ is linearly independent in $P_2(\mathbb{R})$.

\begin{solution}
Set a linear combination equal to the zero polynomial and group by powers of $x$:
\begin{align*}
c_1(1 + x) + c_2(x + x^2) + c_3(1 + x^2) &= 0 + 0x + 0x^2 \\
(c_1 + c_3) + (c_1 + c_2)x + (c_2 + c_3)x^2 &= 0 + 0x + 0x^2.
\end{align*}
For polynomials to be equal, their coefficients must match:
\begin{alignat*}{2}
c_1 + c_3 &= 0 &\qquad& \by{constants} \\
c_1 + c_2 &= 0 && \by{$x$ coefficients} \\
c_2 + c_3 &= 0 && \by{$x^2$ coefficients}
\end{alignat*}
From the first two equations, $c_3 = -c_1$ and $c_2 = -c_1$. Substituting into the third:
\[
-c_1 - c_1 = 0 \implies c_1 = 0,
\]
and then $c_2 = c_3 = 0$. The trivial solution is the only solution, so $S$ is \textbf{linearly independent}.
\end{solution}
\end{example}

\begin{example}[Linear Independence in $M_{2\times 2}(\mathbb{R})$]
Determine whether
\[
S = \left\{ \begin{bmatrix} 1 & 1 \\ 0 & 0 \end{bmatrix}, \begin{bmatrix} 0 & 1 \\ 1 & 0 \end{bmatrix}, \begin{bmatrix} 0 & 0 \\ 1 & 1 \end{bmatrix} \right\}
\]
is linearly independent.

\begin{solution}
Setting a linear combination equal to the zero matrix and adding entry by entry:
\[
c_1 \begin{bmatrix} 1 & 1 \\ 0 & 0 \end{bmatrix} + c_2 \begin{bmatrix} 0 & 1 \\ 1 & 0 \end{bmatrix} + c_3 \begin{bmatrix} 0 & 0 \\ 1 & 1 \end{bmatrix}
= \begin{bmatrix} c_1 & c_1 + c_2 \\ c_2 + c_3 & c_3 \end{bmatrix}
= \begin{bmatrix} 0 & 0 \\ 0 & 0 \end{bmatrix}.
\]
The top-left entry gives $c_1 = 0$ and the bottom-right gives $c_3 = 0$. Then the top-right gives $c_2 = 0$. Only the trivial solution exists, so $S$ is \textbf{linearly independent}.
\end{solution}
\end{example}

% =========================================================================
\subsection{Properties of Dependent Sets}

\begin{theorem}[Dependence Means Redundancy]\label{thm:L2-dep}
A set $S = \{\vec{v}_1, \dots, \vec{v}_k\}$ with $k \ge 2$ is linearly dependent if and only if at least one vector $\vec{v}_i$ is a linear combination of the others.
\end{theorem}

\begin{proof}
\textbf{($\Longrightarrow$)} Suppose $c_1\vec{v}_1 + \dots + c_k\vec{v}_k = \vec{0}$ with some $c_i \neq 0$. Move every other term to the right side and divide by $c_i$:
\[
\vec{v}_i = -\frac{c_1}{c_i}\vec{v}_1 - \dots - \frac{c_{i-1}}{c_i}\vec{v}_{i-1} - \frac{c_{i+1}}{c_i}\vec{v}_{i+1} - \dots - \frac{c_k}{c_i}\vec{v}_k.
\]
So $\vec{v}_i$ is a linear combination of the others.

\textbf{($\Longleftarrow$)} Suppose $\vec{v}_i = \sum_{j \neq i} a_j\vec{v}_j$. Then
\[
a_1\vec{v}_1 + \dots + a_{i-1}\vec{v}_{i-1} + (-1)\vec{v}_i + a_{i+1}\vec{v}_{i+1} + \dots + a_k\vec{v}_k = \vec{0}.
\]
The coefficient of $\vec{v}_i$ is $-1 \neq 0$, so this is a non-trivial solution and $S$ is dependent.
\end{proof}

\begin{theorem}[Quick Facts]
Let $V$ be a vector space.
\begin{enumerate}
    \item Any set containing $\vec{0}$ is linearly dependent.
    \item A single vector $\{\vec{v}\}$ is linearly independent if and only if $\vec{v} \neq \vec{0}$.
    \item Two vectors $\{\vec{v}_1, \vec{v}_2\}$ are linearly dependent if and only if one is a scalar multiple of the other.
\end{enumerate}
\end{theorem}

\begin{proof}
\textbf{(1)} If $\vec{v}_i = \vec{0}$, take $c_i = 1$ and every other coefficient $0$. Then $c_1\vec{v}_1 + \dots + c_k\vec{v}_k = 1 \cdot \vec{0} = \vec{0}$, and this is a non-trivial solution.

\textbf{(2)} If $\vec{v} = \vec{0}$, the set is dependent by (1). If $\vec{v} \neq \vec{0}$ and $c\vec{v} = \vec{0}$, then $c = 0$ by Theorem~\ref{thm:L1-arith}(4), so only the trivial solution exists.

\textbf{(3)} This is Theorem~\ref{thm:L2-dep} with $k = 2$: one vector is a combination of the other exactly when it is a scalar multiple of it.
\end{proof}

\begin{example}[Spotting Dependence by Inspection]
Without row reducing, decide whether each set is linearly independent.
\[
\text{(a)} \ \left\{ \begin{bmatrix} 1 \\ 2 \end{bmatrix}, \begin{bmatrix} -3 \\ -6 \end{bmatrix} \right\}
\qquad
\text{(b)} \ \{x, \ x^2, \ 0\}
\qquad
\text{(c)} \ \{1 + x, \ 1 - x, \ 2\}
\]

\begin{solution}
\textbf{(a)} Dependent: the second vector is $-3$ times the first.

\textbf{(b)} Dependent: the set contains the zero polynomial.

\textbf{(c)} Dependent: adding the first two gives the third,
\[
(1 + x) + (1 - x) - 2 = 0,
\]
which is a non-trivial combination equal to zero.
\end{solution}
\end{example}

% =========================================================================
\subsection{The Importance of Independence}

\begin{theorem}[Unique Representation Theorem]
Let $V$ be a vector space over $\mathbb{R}$ and let $S = \{\vec{v}_1, \dots, \vec{v}_k\}$ be a spanning set for $V$. Then every vector in $V$ can be \emph{uniquely} expressed as a linear combination of the vectors in $S$ if and only if $S$ is linearly independent.
\end{theorem}

\begin{proof}
\textbf{($\Longrightarrow$)} Suppose every vector has a unique representation. The zero vector already has the representation
\[
0\vec{v}_1 + \dots + 0\vec{v}_k = \vec{0}.
\]
If $c_1\vec{v}_1 + \dots + c_k\vec{v}_k = \vec{0}$ as well, then by uniqueness of the representation of $\vec{0}$, we must have $c_1 = \dots = c_k = 0$. Hence $S$ is linearly independent.

\textbf{($\Longleftarrow$)} Suppose $S$ is linearly independent, and let $\vec{v} \in V$. Because $S$ spans $V$, there is at least one representation $\vec{v} = \sum_{i=1}^k c_i\vec{v}_i$. Suppose there is another, $\vec{v} = \sum_{i=1}^k d_i\vec{v}_i$. Subtracting the two:
\[
\sum_{i=1}^k c_i\vec{v}_i - \sum_{i=1}^k d_i\vec{v}_i = \vec{0}
\implies
\sum_{i=1}^k (c_i - d_i)\vec{v}_i = \vec{0}.
\]
Because $S$ is linearly independent, every coefficient must be zero:
\[
c_i - d_i = 0 \implies c_i = d_i \qquad \text{for all } i = 1, \dots, k.
\]
So the two representations are identical, and the linear combination is unique.
\end{proof}

\begin{example}[Unique vs.\ Non-Unique Representations]
Compare how $\mathbb{R}^2$ vectors are built from
\[
S = \left\{ \begin{bmatrix} 1 \\ 1 \end{bmatrix}, \begin{bmatrix} 1 \\ -1 \end{bmatrix} \right\}
\qquad \text{and} \qquad
T = \left\{ \begin{bmatrix} 1 \\ 0 \end{bmatrix}, \begin{bmatrix} 0 \\ 1 \end{bmatrix}, \begin{bmatrix} 1 \\ 1 \end{bmatrix} \right\}.
\]

\begin{solution}
Both sets span $\mathbb{R}^2$. The set $S$ is independent (neither vector is a multiple of the other), so by the earlier spanning example the coefficients are forced:
\[
\begin{bmatrix} 5 \\ 1 \end{bmatrix} = 3\begin{bmatrix} 1 \\ 1 \end{bmatrix} + 2\begin{bmatrix} 1 \\ -1 \end{bmatrix},
\qquad c_1 = \frac{5 + 1}{2} = 3, \quad c_2 = \frac{5 - 1}{2} = 2,
\]
and there is no other choice.

The set $T$ is dependent (the third vector is the sum of the first two), and the same vector can be built in more than one way:
\[
\begin{bmatrix} 2 \\ 3 \end{bmatrix}
= 2\begin{bmatrix} 1 \\ 0 \end{bmatrix} + 3\begin{bmatrix} 0 \\ 1 \end{bmatrix} + 0\begin{bmatrix} 1 \\ 1 \end{bmatrix}
= 1\begin{bmatrix} 1 \\ 0 \end{bmatrix} + 2\begin{bmatrix} 0 \\ 1 \end{bmatrix} + 1\begin{bmatrix} 1 \\ 1 \end{bmatrix}.
\]
This is exactly what the theorem predicts: redundancy in the set means representations stop being unique.
\end{solution}
\end{example}
\newpage
\section{Lecture 3: Bases, Dimension, and Expansion}

\begin{intuition}
A basis is the perfect skeleton of a vector space. It is a set of vectors that spans the entire space (so you can reach absolutely everywhere) but is also linearly independent (so there is absolutely no redundancy in your coordinates). Because of these two properties, a basis provides a perfect, unique coordinate system for any space, ensuring every single vector has exactly one ``address.''
\end{intuition}

% =========================================================================
\subsection{Bases}

\begin{definition}[Basis]
Let $V$ be a vector space. A subset $\mathcal{B} = \{\vec{v}_1, \dots, \vec{v}_k\} \subseteq V$ is called a \textbf{basis} for $V$ if it satisfies two conditions:
\begin{enumerate}
    \item $\mathcal{B}$ is linearly independent (no redundant vectors), and
    \item $\operatorname{Span}(\mathcal{B}) = V$ (enough vectors to build the whole space).
\end{enumerate}
\emph{Convention:} the basis for the trivial subspace $\{\vec{0}\}$ is the empty set $\emptyset$.
\end{definition}

\begin{remark}
By the Unique Representation Theorem from Lecture 2, a basis is exactly a set where every vector of $V$ can be built in \emph{one and only one} way. Spanning gives ``at least one way,'' independence gives ``at most one way.''
\end{remark}

\begin{example}[Basis or Not?]
Decide whether each set is a basis for $\mathbb{R}^2$:
\[
S = \left\{ \begin{bmatrix} 1 \\ 1 \end{bmatrix}, \begin{bmatrix} 1 \\ -1 \end{bmatrix} \right\},
\qquad
T = \left\{ \begin{bmatrix} 1 \\ 0 \end{bmatrix}, \begin{bmatrix} -1 \\ 2 \end{bmatrix}, \begin{bmatrix} 2 \\ 3 \end{bmatrix} \right\}.
\]

\begin{solution}
\textbf{The set $S$.} In Lecture 2 we showed $S$ spans $\mathbb{R}^2$, and neither vector is a multiple of the other, so $S$ is independent. Both conditions hold, so $S$ \textbf{is} a basis.

\textbf{The set $T$.} $T$ spans $\mathbb{R}^2$ (the first two vectors alone already do), but it is dependent:
\[
7\begin{bmatrix} 1 \\ 0 \end{bmatrix} + 3\begin{bmatrix} -1 \\ 2 \end{bmatrix} - 2\begin{bmatrix} 2 \\ 3 \end{bmatrix} = \begin{bmatrix} 7 - 3 - 4 \\ 0 + 6 - 6 \end{bmatrix} = \begin{bmatrix} 0 \\ 0 \end{bmatrix}.
\]
So $T$ is \textbf{not} a basis. Spanning on its own is not enough; there is a redundant vector.
\end{solution}
\end{example}

\subsubsection{Standard Bases}

Just like the $x, y, z$ axes in geometry, our standard vector spaces have ``standard'' bases that are the simplest to work with.

\begin{definition}[Standard Basis]
\begin{itemize}
    \item \textbf{$\mathbb{R}^n$:} let $\vec{e}_i$ be the vector with a $1$ in entry $i$ and $0$ everywhere else. The standard basis is $\{\vec{e}_1, \dots, \vec{e}_n\}$. For example, in $\mathbb{R}^3$:
    \[
    \{\vec{e}_1, \vec{e}_2, \vec{e}_3\} = \left\{ \begin{bmatrix} 1 \\ 0 \\ 0 \end{bmatrix}, \begin{bmatrix} 0 \\ 1 \\ 0 \end{bmatrix}, \begin{bmatrix} 0 \\ 0 \\ 1 \end{bmatrix} \right\}.
    \]
    \item \textbf{$P_n(\mathbb{R})$:} the monomial basis $\{1, x, x^2, \dots, x^n\}$.
    \item \textbf{$M_{m \times n}(\mathbb{R})$:} let $E_{ij}$ be the matrix with a $1$ in position $(i, j)$ and $0$ everywhere else. The standard basis is $\{E_{ij} \mid 1 \le i \le m, \ 1 \le j \le n\}$. For example, in $M_{2 \times 2}(\mathbb{R})$:
    \[
    \{E_{11}, E_{12}, E_{21}, E_{22}\} = \left\{
    \begin{bmatrix} 1 & 0 \\ 0 & 0 \end{bmatrix},
    \begin{bmatrix} 0 & 1 \\ 0 & 0 \end{bmatrix},
    \begin{bmatrix} 0 & 0 \\ 1 & 0 \end{bmatrix},
    \begin{bmatrix} 0 & 0 \\ 0 & 1 \end{bmatrix}
    \right\}.
    \]
\end{itemize}
\end{definition}

\begin{example}[Writing Vectors in a Standard Basis]
Write each vector as a linear combination of the standard basis of its space:
\[
\vec{x} = \begin{bmatrix} 4 \\ -1 \\ 7 \end{bmatrix},
\qquad
p(x) = 4 - x + 2x^3,
\qquad
A = \begin{bmatrix} 3 & -1 \\ 0 & 5 \end{bmatrix}.
\]

\begin{solution}
With the standard basis, the coefficients are just the entries (or coefficients) you can already see:
\begin{align*}
\vec{x} &= 4\vec{e}_1 - 1\vec{e}_2 + 7\vec{e}_3, \\
p(x)    &= 4(1) - 1(x) + 0(x^2) + 2(x^3), \\
A       &= 3E_{11} - 1E_{12} + 0E_{21} + 5E_{22}.
\end{align*}
This is why standard bases are ``standard'': reading off the coefficients takes no work.
\end{solution}
\end{example}

\begin{example}[Verifying That $\{1, x, x^2\}$ Is a Basis for $P_2(\mathbb{R})$]
Prove directly from the definition that $\{1, x, x^2\}$ is a basis for $P_2(\mathbb{R})$.

\begin{solution}
To rigorously prove this is a basis, we must check both criteria in the definition.

\textbf{Spanning.} Let $p(x) \in P_2(\mathbb{R})$ be an arbitrary polynomial. By definition, it has the form $p(x) = a + bx + cx^2$ where $a, b, c \in \mathbb{R}$. We can write this exactly as a linear combination of our basis elements:
\[
p(x) = a(1) + b(x) + c(x^2).
\]
Since any $p(x)$ can be constructed this way, $P_2(\mathbb{R}) = \operatorname{Span}\{1, x, x^2\}$.

\textbf{Linear independence.} Suppose a linear combination equals the zero polynomial, which evaluates to $0$ for every $x \in \mathbb{R}$:
\[
c_1(1) + c_2(x) + c_3(x^2) = 0 \qquad \text{for all } x.
\]
Because this equation is an identity (it must hold for \emph{any} real number $x$), we can strategically plug in values of $x$:
\begin{alignat*}{2}
x = 0: &\quad c_1 = 0 &\qquad& \by{kills the $x$ and $x^2$ terms} \\
x = 1: &\quad c_2 + c_3 = 0 && \\
x = -1: &\quad -c_2 + c_3 = 0 &&
\end{alignat*}
Adding the last two equations gives $2c_3 = 0$, so $c_3 = 0$, and then $c_2 = 0$. The trivial solution is the only solution, so $\{1, x, x^2\}$ is linearly independent.

Because $\{1, x, x^2\}$ spans $P_2(\mathbb{R})$ and is linearly independent, it is a basis for $P_2(\mathbb{R})$.
\end{solution}
\end{example}

% =========================================================================
\subsection{Finding a Basis for a Span}

\begin{theorem}[Removing Redundant Vectors]\label{thm:L3-remove}
Let $\vec{v}_1, \dots, \vec{v}_k \in V$. If a vector $\vec{v}_i$ is a linear combination of the others, it is essentially ``dead weight.'' Removing it does not alter the span:
\[
\operatorname{Span}\{\vec{v}_1, \dots, \vec{v}_k\}
=
\operatorname{Span}\{\vec{v}_1, \dots, \vec{v}_{i-1}, \vec{v}_{i+1}, \dots, \vec{v}_k\}.
\]
\end{theorem}

\begin{proof}
Call the left span $W$ and the right span $W'$. Every combination of the smaller set is also a combination of the bigger set (just give $\vec{v}_i$ the coefficient $0$), so $W' \subseteq W$.

For the other direction, write $\vec{v}_i = \sum_{j \neq i} a_j\vec{v}_j$, and take any $\vec{w} = \sum_{j=1}^k c_j\vec{v}_j \in W$. Substituting for $\vec{v}_i$:
\[
\vec{w} = \sum_{j \neq i} c_j\vec{v}_j + c_i \sum_{j \neq i} a_j\vec{v}_j = \sum_{j \neq i} (c_j + c_i a_j)\vec{v}_j,
\]
which is a combination of the vectors without $\vec{v}_i$. So $W \subseteq W'$, and $W = W'$.
\end{proof}

So to get a basis for $\operatorname{Span}(S)$, we keep throwing away redundant vectors until the set is independent. Row reduction does this in one go.

\begin{formula}[How to Find a Basis for $\operatorname{Span}\{\vec{v}_1, \dots, \vec{v}_k\}$]
\begin{enumerate}
    \item Set up $c_1\vec{v}_1 + \dots + c_k\vec{v}_k = \vec{0}$ and write the coefficient matrix (for vectors in $\mathbb{R}^n$, the columns are just the vectors).
    \item Row reduce to REF (or RREF).
    \item The \textbf{original} vectors that correspond to the pivot columns form a basis. The non-pivot columns are redundant.
\end{enumerate}
\end{formula}

\begin{remark}
Why does this work? Row operations change the rows, but they never change which linear relations hold between the \emph{columns}. So if column 4 of the RREF equals column 1 plus column 3, the same relation holds between the original vectors. Always take the basis vectors from the \emph{original} set, never from the reduced matrix.
\end{remark}

\begin{example}[A Basis for a Span in $\mathbb{R}^4$]
Let
\[
U = \operatorname{Span}\left\{
\begin{bmatrix} 1 \\ 1 \\ -1 \\ 3 \end{bmatrix},
\begin{bmatrix} -2 \\ -2 \\ 2 \\ -6 \end{bmatrix},
\begin{bmatrix} 2 \\ 1 \\ 1 \\ 2 \end{bmatrix},
\begin{bmatrix} 3 \\ 2 \\ 0 \\ 5 \end{bmatrix}
\right\}.
\]
Find a basis for $U$ and $\dim(U)$.

\begin{solution}
To find a basis, we must filter out the redundant vectors. We form the matrix whose columns are the given vectors and row reduce it, which tracks the dependencies between the columns. First we use the leading $1$ in column 1 to clear everything below it. Column 2 then has no pivot, so the next pivot is in column 3:
\begin{align*}
\begin{bmatrix}
 1 & -2 & 2 & 3 \\
 1 & -2 & 1 & 2 \\
-1 &  2 & 1 & 0 \\
 3 & -6 & 2 & 5
\end{bmatrix}
&\xrightarrow[R_4 - 3R_1]{R_2 - R_1,\ R_3 + R_1}
\begin{bmatrix}
1 & -2 &  2 &  3 \\
0 &  0 & -1 & -1 \\
0 &  0 &  3 &  3 \\
0 &  0 & -4 & -4
\end{bmatrix}
\xrightarrow{-R_2}
\begin{bmatrix}
1 & -2 &  2 &  3 \\
0 &  0 &  1 &  1 \\
0 &  0 &  3 &  3 \\
0 &  0 & -4 & -4
\end{bmatrix} \\[0.8em]
&\xrightarrow[R_4 + 4R_2]{R_3 - 3R_2}
\begin{bmatrix}
1 & -2 & 2 & 3 \\
0 &  0 & 1 & 1 \\
0 &  0 & 0 & 0 \\
0 &  0 & 0 & 0
\end{bmatrix}
\xrightarrow{R_1 - 2R_2}
\begin{bmatrix}
1 & -2 & 0 & 1 \\
0 &  0 & 1 & 1 \\
0 &  0 & 0 & 0 \\
0 &  0 & 0 & 0
\end{bmatrix}.
\end{align*}
The pivots are in \textbf{column 1} and \textbf{column 3}, so columns 2 and 4 are redundant. The RREF even tells us how:
\[
\vec{v}_2 = -2\vec{v}_1, \qquad \vec{v}_4 = \vec{v}_1 + \vec{v}_3.
\]
Keeping the \textbf{original vectors} from the pivot columns,
\[
\mathcal{B} = \left\{
\begin{bmatrix} 1 \\ 1 \\ -1 \\ 3 \end{bmatrix},
\begin{bmatrix} 2 \\ 1 \\ 1 \\ 2 \end{bmatrix}
\right\}
\]
is a basis for $U$, and $\dim(U) = 2$.
\end{solution}
\end{example}

\begin{example}[A Basis for a Span in $P_2(\mathbb{R})$]
Let $U = \operatorname{Span}\{2 - x + x^2, \ 1 + x - 2x^2, \ 3 - x^2\}$. Find a basis for $U$.

\begin{solution}
Suppose
\[
c_1(2 - x + x^2) + c_2(1 + x - 2x^2) + c_3(3 - x^2) = 0.
\]
Expanding and grouping by powers of $x$, each coefficient must be $0$:
\begin{alignat*}{2}
2c_1 + c_2 + 3c_3 &= 0 &\qquad& \by{constants} \\
-c_1 + c_2 &= 0 && \by{$x$ terms} \\
c_1 - 2c_2 - c_3 &= 0 && \by{$x^2$ terms}
\end{alignat*}
Row reducing the coefficient matrix:
\begin{align*}
\begin{bmatrix}
 2 &  1 &  3 \\
-1 &  1 &  0 \\
 1 & -2 & -1
\end{bmatrix}
&\xrightarrow{R_1 \leftrightarrow R_3}
\begin{bmatrix}
 1 & -2 & -1 \\
-1 &  1 &  0 \\
 2 &  1 &  3
\end{bmatrix}
\xrightarrow[R_3 - 2R_1]{R_2 + R_1}
\begin{bmatrix}
 1 & -2 & -1 \\
 0 & -1 & -1 \\
 0 &  5 &  5
\end{bmatrix} \\[0.8em]
&\xrightarrow{-R_2}
\begin{bmatrix}
 1 & -2 & -1 \\
 0 &  1 &  1 \\
 0 &  5 &  5
\end{bmatrix}
\xrightarrow[R_1 + 2R_2]{R_3 - 5R_2}
\begin{bmatrix}
 1 & 0 & 1 \\
 0 & 1 & 1 \\
 0 & 0 & 0
\end{bmatrix}.
\end{align*}
The pivots appear in the first and second columns, so the third polynomial is redundant. Indeed, the third column says it is the sum of the first two:
\[
(2 - x + x^2) + (1 + x - 2x^2) = 3 - x^2.
\]
Keeping the first two original polynomials,
\[
\mathcal{B} = \{2 - x + x^2, \ 1 + x - 2x^2\}
\]
is a basis for $U$, and $\dim(U) = 2$.
\end{solution}
\end{example}

% =========================================================================
\subsection{Dimension of a Vector Space}

\begin{theorem}[Dimension Is Well Defined]
If $\mathcal{B} = \{\vec{v}_1, \dots, \vec{v}_k\}$ and $\mathcal{C} = \{\vec{w}_1, \dots, \vec{w}_\ell\}$ are both bases for a vector space $V$, then $k = \ell$.
\end{theorem}

\begin{remark}
This is a mathematical guarantee: a vector space always requires the exact same number of basis vectors, regardless of which specific vectors you choose. The idea behind the proof is that if $k < \ell$, then $\ell$ vectors would be sitting inside the span of only $k$ vectors, which forces them to be dependent. So $\mathcal{C}$ could not be a basis.
\end{remark}

\begin{definition}[Dimension]
If $\mathcal{B} = \{\vec{v}_1, \dots, \vec{v}_k\}$ is a basis for a vector space $V$, then the \textbf{dimension} of $V$ is the number of vectors in $\mathcal{B}$:
\[
\dim(V) = k.
\]
In particular, if $V = \{\vec{0}\}$, then $\dim(V) = 0$, since its basis is $\emptyset$. If $V$ has a finite basis, we call $V$ \textbf{finite dimensional}.
\end{definition}

\begin{example}[Dimensions of the Standard Spaces]
Find $\dim(\mathbb{R}^n)$, $\dim(P_n(\mathbb{R}))$ and $\dim(M_{m \times n}(\mathbb{R}))$.

\begin{solution}
Count the vectors in each standard basis:
\begin{alignat*}{2}
\dim(\mathbb{R}^n) &= n &\qquad& \by{one $\vec{e}_i$ per entry} \\
\dim(P_n(\mathbb{R})) &= n + 1 && \by{$1, x, \dots, x^n$; the $+1$ is for the constant term} \\
\dim(M_{m \times n}(\mathbb{R})) &= mn && \by{one $E_{ij}$ per matrix entry}
\end{alignat*}
For example, $\dim(P_2(\mathbb{R})) = 3$ and $\dim(M_{2 \times 3}(\mathbb{R})) = 6$.
\end{solution}
\end{example}

\begin{example}[Dimension of a Subspace Given by a Condition]
Find a basis and the dimension of
\[
S = \left\{ \begin{bmatrix} x_1 \\ x_2 \\ x_3 \\ x_4 \end{bmatrix} \in \mathbb{R}^4 \;\middle|\; x_1 + x_2 + x_3 + x_4 = 0 \right\}.
\]

\begin{solution}
Solve the condition for one variable, $x_4 = -x_1 - x_2 - x_3$, and split the general vector by its free variables:
\[
\begin{bmatrix} x_1 \\ x_2 \\ x_3 \\ -x_1 - x_2 - x_3 \end{bmatrix}
= x_1 \begin{bmatrix} 1 \\ 0 \\ 0 \\ -1 \end{bmatrix}
+ x_2 \begin{bmatrix} 0 \\ 1 \\ 0 \\ -1 \end{bmatrix}
+ x_3 \begin{bmatrix} 0 \\ 0 \\ 1 \\ -1 \end{bmatrix}.
\]
So these three vectors span $S$. They are also independent: if a combination of them is $\vec{0}$, the first three entries give $x_1 = x_2 = x_3 = 0$ directly. So they form a basis, and
\[
\dim(S) = 3.
\]
This makes sense: one equation in $\mathbb{R}^4$ removes one degree of freedom, $4 - 1 = 3$.
\end{solution}
\end{example}

\begin{example}[Dimension of the Symmetric $2 \times 2$ Matrices]
Find a basis and the dimension of $S = \{A \in M_{2 \times 2}(\mathbb{R}) \mid A = A^T\}$.

\begin{solution}
Write a general $2 \times 2$ matrix with entries $a, b, c, d$. The condition $A = A^T$ means
\[
\begin{bmatrix} a & b \\ c & d \end{bmatrix} = \begin{bmatrix} a & c \\ b & d \end{bmatrix}
\implies b = c.
\]
So every element of $S$ looks like
\[
\begin{bmatrix} a & b \\ b & d \end{bmatrix}
= a\begin{bmatrix} 1 & 0 \\ 0 & 0 \end{bmatrix}
+ b\begin{bmatrix} 0 & 1 \\ 1 & 0 \end{bmatrix}
+ d\begin{bmatrix} 0 & 0 \\ 0 & 1 \end{bmatrix}.
\]
These three matrices span $S$, and they are independent (if the combination is the zero matrix, then $a = b = d = 0$ by reading off entries). So they form a basis, and $\dim(S) = 3$.
\end{solution}
\end{example}

\subsubsection{Counting Vectors}

\begin{theorem}[Too Many, Too Few, Just Right]\label{thm:L3-count}
Let $V$ be a finite dimensional vector space with $\dim(V) = n$. Then:
\begin{enumerate}
    \item A set of \textbf{more than $n$} vectors must be linearly dependent. (You have too many vectors, creating guaranteed redundancy.)
    \item A set of \textbf{fewer than $n$} vectors cannot span $V$. (You don't have enough directions to reach every corner of the space.)
    \item A set of \textbf{exactly $n$} vectors spans $V$ if and only if it is linearly independent. (If the count is exactly $n$, passing one test automatically guarantees passing the other.)
\end{enumerate}
\end{theorem}

\begin{formula}[Basis Check When $\dim(V) = n$ Is Known]
Let $\mathcal{B}$ have $k$ vectors.
\begin{enumerate}
    \item If $k \neq n$, then $\mathcal{B}$ is \textbf{not} a basis. Done.
    \item If $k = n$, check only linear independence. Independent means basis, dependent means not a basis.
\end{enumerate}
\end{formula}

\begin{example}[Counting Alone Decides]
Decide whether each set is a basis.
\[
\text{(a)} \ \{1, \ x, \ x^2, \ 2x^2, \ 1 - 3x\} \subseteq P_2(\mathbb{R})
\qquad
\text{(b)} \ \left\{ \begin{bmatrix} 1 \\ 0 \\ 0 \end{bmatrix}, \begin{bmatrix} 0 \\ 1 \\ 0 \end{bmatrix} \right\} \subseteq \mathbb{R}^3
\]

\begin{solution}
\textbf{(a)} The set has 5 elements, but $\dim(P_2(\mathbb{R})) = 3$. Because $5 > 3$, there is unavoidable redundancy, so the set is dependent and \textbf{not} a basis.

\textbf{(b)} There are only 2 vectors and $\dim(\mathbb{R}^3) = 3$, so this set cannot span $\mathbb{R}^3$. It only describes a 2D plane (the $xy$-plane), so it is \textbf{not} a basis.
\end{solution}
\end{example}

\begin{example}[Right Count, Check Independence]
Is $\mathcal{B} = \{x + x^2, \ 2 - x^2, \ 1 + x + x^2\}$ a basis for $P_2(\mathbb{R})$?

\begin{solution}
$\mathcal{B}$ has 3 vectors and $\dim(P_2(\mathbb{R})) = 3$, so we only need independence. Setting
\[
c_1(x + x^2) + c_2(2 - x^2) + c_3(1 + x + x^2) = 0
\]
and matching coefficients:
\begin{alignat*}{2}
2c_2 + c_3 &= 0 &\qquad& \by{constants} \\
c_1 + c_3 &= 0 && \by{$x$ terms} \\
c_1 - c_2 + c_3 &= 0 && \by{$x^2$ terms}
\end{alignat*}
Subtracting the second equation from the third gives $-c_2 = 0$, so $c_2 = 0$. Then the first gives $c_3 = 0$, and the second gives $c_1 = 0$. Only the trivial solution exists, so $\mathcal{B}$ is independent, and therefore a basis for $P_2(\mathbb{R})$. We never had to check spanning.
\end{solution}
\end{example}

\begin{example}[Right Count, Still Fails]
Is
\[
\mathcal{B} = \left\{ \begin{bmatrix} 1 \\ -1 \\ 0 \end{bmatrix}, \begin{bmatrix} 2 \\ 1 \\ -1 \end{bmatrix}, \begin{bmatrix} -3 \\ 0 \\ 1 \end{bmatrix} \right\}
\]
a basis for $\mathbb{R}^3$?

\begin{solution}
The count is right ($3 = \dim(\mathbb{R}^3)$), so we check independence. Adding all three vectors:
\[
\begin{bmatrix} 1 \\ -1 \\ 0 \end{bmatrix} + \begin{bmatrix} 2 \\ 1 \\ -1 \end{bmatrix} + \begin{bmatrix} -3 \\ 0 \\ 1 \end{bmatrix} = \begin{bmatrix} 0 \\ 0 \\ 0 \end{bmatrix}.
\]
This is a non-trivial combination equal to zero, so $\mathcal{B}$ is dependent and \textbf{not} a basis.
\end{solution}
\end{example}

% =========================================================================
\subsection{Basis Expansion}

\begin{lemma}[Adding a Vector Outside the Span]\label{lem:L3-add}
Let $C = \{\vec{v}_1, \dots, \vec{v}_k\}$ be linearly independent. If $\vec{v} \notin \operatorname{Span}(C)$, then $\{\vec{v}_1, \dots, \vec{v}_k, \vec{v}\}$ is still linearly independent.
\end{lemma}

\begin{proof}
Suppose $c_1\vec{v}_1 + \dots + c_k\vec{v}_k + a\vec{v} = \vec{0}$. If $a \neq 0$, we could solve
\[
\vec{v} = -\frac{c_1}{a}\vec{v}_1 - \dots - \frac{c_k}{a}\vec{v}_k,
\]
which would put $\vec{v}$ in $\operatorname{Span}(C)$, a contradiction. So $a = 0$, and the equation becomes $c_1\vec{v}_1 + \dots + c_k\vec{v}_k = \vec{0}$. Since $C$ is independent, all $c_i = 0$. Only the trivial solution exists.
\end{proof}

\begin{theorem}[Basis Expansion]
Let $V$ be an $n$-dimensional vector space. If $C = \{\vec{v}_1, \dots, \vec{v}_k\} \subseteq V$ is linearly independent with $k < n$, then we can find additional vectors $\vec{v}_{k+1}, \dots, \vec{v}_n \in V$ such that
\[
\{\vec{v}_1, \dots, \vec{v}_k, \vec{v}_{k+1}, \dots, \vec{v}_n\}
\]
is a basis for $V$. In other words, you can always ``pad'' an independent set until it becomes a basis.
\end{theorem}

\begin{proof}
Since $k < n$, $C$ cannot span $V$ (Theorem~\ref{thm:L3-count}(2)), so there is some $\vec{v}_{k+1} \in V$ with $\vec{v}_{k+1} \notin \operatorname{Span}(C)$. By Lemma~\ref{lem:L3-add}, $C \cup \{\vec{v}_{k+1}\}$ is still independent. Repeat this until the set has $n$ vectors. An independent set of exactly $n$ vectors is a basis by Theorem~\ref{thm:L3-count}(3).
\end{proof}

\begin{formula}[How to Expand to a Basis]
\begin{enumerate}
    \item Start with the independent set $C$ and note how many vectors are missing ($n - k$).
    \item Set a generic vector equal to a combination of $C$ and row reduce to find the condition for it to be \emph{inside} $\operatorname{Span}(C)$.
    \item Choose a new vector that intentionally violates that condition, and add it. Repeat if more vectors are needed.
\end{enumerate}
\end{formula}

\begin{example}[Expanding in $\mathbb{R}^3$]
Extend
\[
C = \left\{ \begin{bmatrix} 1 \\ 2 \\ 3 \end{bmatrix}, \begin{bmatrix} 0 \\ 1 \\ 0 \end{bmatrix} \right\}
\]
to a basis for $\mathbb{R}^3$.

\begin{solution}
$C$ is linearly independent (the vectors are not scalar multiples). Since $\dim(\mathbb{R}^3) = 3$, we need exactly one more vector, and it must \textbf{not} be a linear combination of the vectors in $C$.

Take a generic vector with entries $a, b, c$ and ask when there exist $c_1, c_2$ with
\[
c_1 \begin{bmatrix} 1 \\ 2 \\ 3 \end{bmatrix} + c_2 \begin{bmatrix} 0 \\ 1 \\ 0 \end{bmatrix} = \begin{bmatrix} a \\ b \\ c \end{bmatrix}.
\]
Row reducing the augmented matrix:
\[
\left[\begin{array}{cc|c}
1 & 0 & a \\
2 & 1 & b \\
3 & 0 & c
\end{array}\right]
\xrightarrow[R_3 - 3R_1]{R_2 - 2R_1}
\left[\begin{array}{cc|c}
1 & 0 & a \\
0 & 1 & b - 2a \\
0 & 0 & c - 3a
\end{array}\right].
\]
Looking at the bottom row, the vector is inside the span exactly when $c - 3a = 0$. Choosing any $a, b, c$ that deliberately break this rule gives a valid extension vector.

Take $a = 2$, $b = 0$, $c = 1$. Since $1 \neq 3(2)$, this vector is outside $\operatorname{Span}(C)$, so by Lemma~\ref{lem:L3-add} the expanded set
\[
\mathcal{B} = \left\{ \begin{bmatrix} 1 \\ 2 \\ 3 \end{bmatrix}, \begin{bmatrix} 0 \\ 1 \\ 0 \end{bmatrix}, \begin{bmatrix} 2 \\ 0 \\ 1 \end{bmatrix} \right\}
\]
is independent. It has $3 = \dim(\mathbb{R}^3)$ vectors, so it is a full basis for $\mathbb{R}^3$.
\end{solution}
\end{example}

\begin{example}[Expanding in $P_3(\mathbb{R})$]
Extend $C = \{x, \ 1 + x\}$ to a basis for $P_3(\mathbb{R})$.

\begin{solution}
The space $P_3(\mathbb{R})$ is 4-dimensional, so we need two more polynomials. Notice that $C$ contains only polynomials of degree at most 1. Any combination of them has degree at most 1, so it can never produce $x^2$:
\[
x^2 \notin \operatorname{Span}\{x, \ 1 + x\}.
\]
By Lemma~\ref{lem:L3-add}, $\{x, \ 1 + x, \ x^2\}$ is independent. By the same reasoning, nothing in its span has an $x^3$ term, so $x^3$ can be added too. This gives
\[
\mathcal{B} = \{x, \ 1 + x, \ x^2, \ x^3\},
\]
an independent set of $4 = \dim(P_3(\mathbb{R}))$ polynomials, which is a basis for $P_3(\mathbb{R})$.
\end{solution}
\end{example}

\end{document}